\section*{الفصل \ref{ch:g3:separating-mixtures} --- فصل مكونات الخلائط}

\begin{solution}{exo:g3:separating-mixtures:1}
\omterm{def:g3:separating-mixtures:sieve}{الغربلة}: حبات الأرز أصغر من حبات البازلاء فتسقط عبر غربال ثقوبه
مناسبة. (والفرز باليد ينجح أيضًا، لكن ببطء.)
\end{solution}

\begin{solution}{exo:g3:separating-mixtures:2}
ينزل التراب ببطء ويكوّن طبقة في القاع؛ ويصير الماء فوقه أصفى. وهذا
هو الترسيب.
\end{solution}

\begin{solution}{exo:g3:separating-mixtures:3}
\omterm{def:g3:separating-mixtures:filter}{الراشح} هو السائل الذي يمر عبر المرشِّح. وحين نحضّر القهوة يكون \omterm{def:g3:separating-mixtures:filter}{الراشح} هو
القهوة في الإبريق؛ ويبقى البن المطحون في
الورقة.
\end{solution}

\begin{solution}{exo:g3:separating-mixtures:4}
نعم. فالملح ذائب، والمرشِّح لا يحجز ما هو ذائب: فهو يمر مع
الماء.
\end{solution}

\begin{solution}{exo:g3:separating-mixtures:5}
\omterm{def:g3:separating-mixtures:evaporate}{بتبخير} الماء: في الشمس، أو بتسخين هادئ، يذهب الماء إلى الهواء ويبقى
السكر في الصحن.
\end{solution}

\begin{solution}{exo:g3:separating-mixtures:6}
يبقى الحصى في الغربال. أما حبيبات الرمل فأصغر من الثقوب، ولذلك
تسقط عبرها.
\end{solution}

\begin{solution}{exo:g3:separating-mixtures:7}
(أ) \omterm{def:g3:separating-mixtures:sieve}{الغربلة}. (ب) \omterm{def:g3:separating-mixtures:filter}{الترشيح} (أو الترسيب ثم \omterm{def:g3:separating-mixtures:decant}{الترويق}). (ج) \omterm{def:g3:separating-mixtures:evaporate}{التبخير}.
\end{solution}

\begin{solution}{exo:g3:separating-mixtures:8}
أربع مراحل: المصافي، وأحواض الترسيب، ومرشِّحات الرمل، وقتل الجراثيم.
وأحواض الترسيب هي التي تزيل الطين.
\end{solution}

\begin{solution}{exo:g3:separating-mixtures:9}
رشّح؛ اجمع الرمل من على الورقة؛ بخّر \omterm{def:g3:separating-mixtures:filter}{الراشح} لتحصل على
الملح.
\end{solution}

\begin{solution}{exo:g3:separating-mixtures:10}
لا. فملح ماء البحر ذائب، والمرشِّح لا يحجز ما هو ذائب: فيبقى \omterm{def:g3:separating-mixtures:filter}{الراشح}
ماءً مالحًا.
\end{solution}

\begin{solution}{exo:g3:separating-mixtures:11}
نصب الشاي برفق بعد أن تترسب الأوراق (\omterm{def:g3:separating-mixtures:decant}{الترويق})، أو نصبه عبر مصفاة أو
مرشِّح (\omterm{def:g3:separating-mixtures:sieve}{الغربلة} أو \omterm{def:g3:separating-mixtures:filter}{الترشيح}).
\end{solution}
